\chapter{Exercise groups}
\label{ch:groups}

Exercise groups get one private room each. Only the staff and the members of that group are
in it, so it is the place for group questions and for feedback on exercise sheets. The usual
order is \textbf{create} \then \textbf{assign} \then \textbf{post}.
All three commands go through a plan that you confirm with \cmd{!confirm}.

\section{Create group rooms}

\begin{you}
!groups create CS101 8
\end{you}

This creates the rooms \enquote{Group 1} to \enquote{Group 8} in the course. Group rooms are
encrypted and invite-only: the owner and the TAs are invited right away, students only
through \cmd{!groups assign}. If the course has groups already, the numbering continues
(\cmd{!groups create CS101 2} then adds Group 9 and Group 10). You can create up to 50
groups at a time.

\section{Put people into groups}

\subsection*{Automatically}

\begin{you}
!groups assign CS101 --auto
\end{you}

Splits all students who have \emph{joined} the course space evenly across the existing
groups. People who haven't accepted their invitation yet are not included, so wait until
most students are in.

\subsection*{With a list}

Send \cmd{!groups assign CS101} and then a CSV file:

\begin{codefile}[groups.csv]
email,group,role
anna.beispiel@stud.unibas.ch,1,
max.muster@stud.unibas.ch,1,
eva.example@stud.unibas.ch,2,
ta.one@unibas.ch,1,ta
\end{codefile}

\begin{itemize}
  \item \cmd{email}: the person's university address. A column \cmd{user} with Matrix IDs
        works too.
  \item \cmd{group}: the group number. The group room must exist already.
  \item \cmd{role} (optional): \cmd{ta} makes the person moderator in that group room, e.g.\
        the TA who runs the group. Empty or \cmd{student} for everyone else.
\end{itemize}

Everyone assigned gets an invitation to their group room. Rows the bot can't use (unknown
group, not a university address) are listed in the plan; the rest is assigned.

\begin{tip}[Excel]
Save the table with \ui{File} \then \ui{Save As} \then \ui{CSV UTF-8}.
Both \cmd{,} and \cmd{;} work as separators, and the first row must contain the column names.
\end{tip}

\section{Post a message per group or per student}

Typical use: points or feedback for an exercise sheet. Send \cmd{!groups post CS101} and then
a CSV file in one of two forms:

\begin{codefile}[one message per group room]
group,message
1,"Sheet 3: 18/20 points, well done"
2,"Sheet 3: 15/20 points, see the comments on task 2"
\end{codefile}

\begin{codefile}[one direct message per student]
email,message
anna.beispiel@stud.unibas.ch,"Sheet 3: 9/10 points"
max.muster@stud.unibas.ch,"Sheet 3: 7/10 points"
\end{codefile}

Put messages that contain commas in quotes. With \cmd{group}, the bot posts the message in
that group's room. With \cmd{email}, it sends the message to the student in a direct chat with
the bot.

The bot always shows a preview of the first messages and waits for \cmd{!confirm}. Once the
messages are sent, they are deleted from the bot's database, because they may contain grades.

\begin{warning}[Check before you confirm]
Posted messages can't be taken back by the bot. A wrong message in a group room can be
deleted in Element (staff are moderators there). Messages sent to a student directly can't
be deleted; send a correction the same way.
\end{warning}
